\chapter{Building Hosts and Clients}
\label{ch:building-hosts}

The host is the only component that can see all of it. The model proposes, the
server executes, the user watches, and when something goes wrong the host was
the only party in a position to have prevented it. That makes it the accountable
component whether or not anybody designed it to be, which is worth deciding
deliberately rather than discovering during an incident.

Building a server is mostly craft: implement the methods, honour the contract,
test the wire. Building a host is mostly judgement. This chapter is about the
five judgements that matter, two of which get skipped in every prototype and
retrofitted in almost none.

\section{What a host is actually for}

A client is thin: attach \texttt{\_meta}, send, parse, run the MRTR loop. A few
hundred lines, and Meridian's is one file.

A host is the hard part, and it has five jobs:

\begin{enumerate}[leftmargin=1.6em, itemsep=3pt]
\item Manage many server connections, with discovery, caching, and reconnection.
\item Merge their catalogues into one namespace the model can reason about.
\item Route the model's intents back to the right server.
\item Enforce a context budget across all of them.
\item Be the consent, audit, and policy boundary.
\end{enumerate}

Jobs four and five are the ones that get skipped in a prototype and then never
retrofitted.

\bookfig[1.0]{host-anatomy}{The host is the only component that sees both
sides. That is why consent, audit, and the context budget can live nowhere
else.}{host-anatomy}

\section{Connecting}

\begin{py}
def connect(self, label: str, transport, *, input_provider=None,
            auth_context: str = "anon") -> ServerBinding:
    client = Client(
        transport,
        name=self.name, version=self.version,
        capabilities=self.capabilities,
        cache=self.cache,
        input_provider=input_provider,
        server_label=label,
        auth_context=auth_context,
    )
    binding = ServerBinding(label=label, client=client,
                            connected_at=time.time())
    with self._lock:
        self.bindings[label] = binding
    return binding
\end{py}

One client per server, as the architecture requires. The label is the namespace
prefix, and it must be stable across restarts because it ends up in
conversation history and audit logs.

\subsection{Discover in parallel, or pay the sum}

The first mistake almost everybody ships:

\begin{py}
def discover_all(self, parallel: bool = True) -> dict[str, dict]:
    """Fetch every server's capabilities.

    In parallel, because these are independent and doing them in sequence
    is the first unnecessary serialisation most hosts ship with.
    """
    def one(binding):
        try:
            result = binding.client.discover()
            binding.capabilities = result.get("capabilities", {})
            binding.healthy = True
            return binding.label, result
        except Exception as exc:
            binding.healthy = False
            binding.last_error = f"{type(exc).__name__}: {exc}"
            return binding.label, {}
    ...
\end{py}

Twelve servers at 60\,ms each is 720\,ms serially and 60\,ms in parallel. On a
desktop application that is the difference between an app that feels ready and
one that shows a spinner every launch.

Note the exception handling. \textbf{One unreachable server must not prevent the
other eleven from working.} A host that fails startup because a rarely used
server is down has coupled everything to everything.

\begin{py}
def test_discovery_reports_capabilities_per_server(self):
    discovered = self.host.discover_all()
    self.assertIn("tools", discovered["risk"]["capabilities"])
    self.assertIn(tasks_ext.EXTENSION_ID,
                  discovered["risk"]["capabilities"]["extensions"])
\end{py}

\subsection{Honour the discovery TTL}

\texttt{server/discover} is cacheable, and Meridian's servers say one hour.
Re-discovering on every task wastes a round trip per server per task for
information that changes when somebody deploys.

The shared \texttt{ResultCache} handles it, which means the host gets this for
free by not doing anything special.

\section{Namespacing, which is not optional}

Tool names are unique \emph{per server}. Two servers each exposing
\texttt{search} is normal and expected. Flatten them into one catalogue and one
of them disappears.

\begin{py}
def catalogue(self) -> list[dict]:
    """The merged, namespaced, budgeted tool list the model actually sees.

    Namespacing is not optional once you have more than one server. Tool
    names are unique per server and nothing more, so two servers each
    offering `search` is normal, expected, and fatal if you flatten them.
    """
    out = []
    with self._lock:
        for label in sorted(self.bindings):
            binding = self.bindings[label]
            for tool in binding.tools:
                entry = dict(tool)
                entry["name"] = f"{label}{NAMESPACE_SEPARATOR}{tool['name']}"
                entry["_server"] = label
                out.append(entry)
    ...
\end{py}

\begin{py}
def test_namespacing_keeps_collisions_apart(self):
    """Two servers may each define the same tool name. That is legal."""
    host = Host()
    host.connect("alpha", InProcessTransport(make("alpha")))
    host.connect("beta", InProcessTransport(make("beta")))
    host.refresh_catalogue()

    names = [t["name"] for t in host.catalogue()]
    self.assertEqual(sorted(names), ["alpha.search", "beta.search"])
    self.assertEqual(
        host.call_tool("alpha.search")["content"][0]["text"], "from alpha")
    self.assertEqual(
        host.call_tool("beta.search")["content"][0]["text"], "from beta")
\end{py}

Note \texttt{sorted(self.bindings)}. Deterministic ordering across servers, for
the same prompt-cache reason as Chapter~\ref{ch:tools}. A host that iterates a
dictionary in insertion order will reorder its catalogue whenever a server
reconnects, and quietly lose every prompt-cache hit.

\begin{notebox}{Choosing a separator}
Meridian uses \texttt{.}, giving \texttt{risk.assess\_account\_risk}. Tool names
allow letters, digits, underscore, hyphen, and dot, so any of them works.

Whatever you pick, pick it once. The namespaced name appears in conversation
history and in audit logs, and changing the separator invalidates both.
\end{notebox}

\section{The context budget}

\keyidea{The host is the context window's memory allocator. Nothing else can be:
a server cannot see the other servers, and the model cannot see its own budget.
If the host does not enforce a limit, there is no limit.}

Connect to eight servers with twenty tools each and your model spends 30,000
tokens per turn on descriptions before it reads a word of the actual task.

\begin{py}
def _apply_budget(self, tools: list[dict], budget: int) -> list[dict]:
    """Drop tools once the catalogue exceeds its token allowance.

    Crude on purpose. The interesting question is not the eviction policy,
    it is that a budget exists at all: without one, a host wired to eight
    servers silently spends thirty thousand tokens per turn on descriptions
    of tools the model will never call.
    """
    kept, spent = [], 0
    for tool in tools:
        cost = estimate_tokens(
            f"{tool.get('name','')}{tool.get('description','')}"
            f"{tool.get('inputSchema','')}")
        if spent + cost > budget:
            continue
        kept.append(tool)
        spent += cost
    return kept
\end{py}

Truncation is the crudest possible policy and I am not defending it. Better ones,
roughly in order of effort:

\begin{itemize}[leftmargin=1.6em, itemsep=3pt]
\item \textbf{Relevance filtering.} Retrieve the tools plausibly related to the
current task and send only those. Costs a retrieval step, saves thousands of
tokens.
\item \textbf{Tiering.} Always include a small core; include the rest only when
the task looks like it needs them.
\item \textbf{Progressive disclosure.} Expose a \texttt{list\_more\_tools} tool.
The model asks when it needs to.
\item \textbf{Just delete some.} Genuinely the best option, and the one people
skip because it requires a conversation with a team.
\end{itemize}

Whatever you choose, measure it:

\begin{py}
def catalogue_tokens(self) -> int:
    return sum(
        estimate_tokens(f"{t.get('name','')}{t.get('description','')}"
                        f"{t.get('inputSchema','')}{t.get('outputSchema','')}")
        for t in self.catalogue())
\end{py}

\begin{measurebox}{Meridian's merged catalogue}
\begin{tabular}{@{}lrr@{}}
\toprule
\thead{Configuration} & \thead{Tools} & \thead{Tokens per turn} \\
\midrule
Four servers, task-shaped catalogue & 10 & 1{,}246 \\
Four servers, REST-mirror catalogue & 38 & 5{,}617 \\
\bottomrule
\end{tabular}

\vspace{6pt}
\footnotesize
This is the number to put on a dashboard. It moves whenever anybody adds a tool
anywhere, and nobody notices until it is 12,000.
\end{measurebox}

\section{Routing}

Split the namespaced name, find the binding, apply consent, call.

\begin{py}
def call_tool(self, qualified: str, arguments: dict | None = None, **kw) -> dict:
    """Route one call, after the consent gate."""
    label, name = self.split_name(qualified)
    binding = self.bindings.get(label)
    if binding is None:
        raise McpError(-32602, f"No server named {label!r}")

    tool = self.find_tool(qualified) or {}
    decision = self.consent.check(qualified, tool, arguments or {})
    self.audit.append({
        "tool": qualified, "decision": decision, "at": time.time(),
        "arguments": arguments or {},
    })
    if decision == ConsentDecision.DENY:
        return tool_error(f"The user declined the call to {qualified}.")

    return binding.client.call_tool(name, arguments or {}, **kw)
\end{py}

Two details worth copying.

\textbf{A denial returns a tool execution error, not an exception.} The model should learn that the call was refused and adapt. ``The user declined'' is information the model can use.

\textbf{Audit before executing.} The record is written whatever happens next,
including if the server times out. An audit log that only records successes is
not an audit log.

\subsection{Parallel fan-out}

\begin{py}
with concurrent.futures.ThreadPoolExecutor(max_workers=len(calls)) as pool:
    futures = [pool.submit(self.call_tool, c["name"], c.get("arguments"))
               for c in calls]
    results = []
    for future in futures:
        try:
            results.append(future.result())
        except McpError as exc:
            results.append(tool_error(f"{exc.message}"))
    return results
\end{py}

Iterating futures in submission order preserves result ordering regardless of
completion order, and catching per-future means one failure produces one error
result rather than discarding two successful calls you already paid for.

\section{Consent that people do not rage-click}

Prompt for everything and users learn to click Approve without reading, which is
worse than not prompting because now you have the audit trail of consent without
any of the reality. Prompt for nothing and the model gets to do anything.

Meridian's default policy:

\begin{py}
class ConsentPolicy:
    """Where the human stays in the loop.

    The default says yes to anything a server marked read-only and asks about
    everything else. That is a defensible default, and it is also exactly the
    kind of thing an attacker will try to talk their way past, which is why
    Chapter 19 argues the annotations must be treated as untrusted unless the
    server is.
    """
\end{py}

Four rules I would defend in a design review:

\textbf{Auto-approve read-only tools from trusted servers.} Reading is
reversible. The word ``trusted'' is load-bearing, and Chapter~\ref{ch:threats}
explains what happens when it is misplaced.

\textbf{Show the arguments, not the tool name.} ``Call \texttt{send\_email}?'' is
not consent. ``Send email to \texttt{board@example.com} with subject
\texttt{Q3 numbers}?'' is. The dangerous part is always in the arguments.

\textbf{Let users grant scoped standing approval.} ``Always allow
\texttt{risk.assess\_account\_risk}'' is reasonable. ``Always allow everything''
is a checkbox that ends careers.

\textbf{Batch related prompts.} A model fanning out to three read tools should
produce one prompt, not three. Three prompts train the click reflex.

\begin{py}
def test_read_only_hint_is_auto_allowed(self):
    seen = []
    host = wire_host(consent=ConsentPolicy(
        auto_allow_read_only=True,
        on_ask=lambda name, args: seen.append(name) or True))
    host.call_tool("risk.assess_account_risk", {"accountId": "ACC-1000"})
    self.assertEqual(seen, [], "a read-only tool should not have prompted")
\end{py}

\section{Handling the three interruptions}

A host loop has to cope with three things that are not ``here is your answer'',
and none of them should block the others.

\subsection{\texttt{input\_required}}

The MRTR loop lives in the client, so the host's involvement is supplying an
\texttt{InputProvider}:

\begin{py}
class InputProvider:
    """How a host answers a server's questions.

    A real host renders a form and waits for a human. Tests supply canned
    answers. Both satisfy this interface, which is the point: the MRTR loop
    below does not care which one it is talking to.
    """
    def elicit(self, key: str, params: dict) -> dict:
        return {"action": "decline"}
\end{py}

The default declines. That is deliberate: a host that has not been wired up should refuse. Inventing answers on the
user's behalf is worse than declining.

\subsection{Task updates}

Polling must not block the loop. Poll on a background thread, surface status to
the user, and let the model keep working on anything that does not depend on the
task.

\subsection{Subscription events}

A \texttt{subscriptions/listen} stream delivers notifications whenever the server
feels like it, which is to say in the middle of something else. Route them to
cache invalidation and carry on:

\begin{py}
def on_notification(self, msg: dict) -> None:
    method = msg.get("method", "")
    if method == "notifications/tools/list_changed":
        self.cache.invalidate_method(self.label, "tools/list")
    elif method == "notifications/resources/updated":
        uri = (msg.get("params") or {}).get("uri")
        if uri:
            self.cache.invalidate_uri(self.label, uri)
\end{py}

A tools-list change mid-task deserves a thought. Refresh immediately and the
model may see the catalogue change between turns, which is confusing but
correct. Defer until the task ends and the model may call a tool that no longer
exists, which produces a clean error it can recover from. Meridian defers, and I
think that is right: catalogue churn mid-task is rarer than task coherence
matters.

\section{Talking to servers from both eras}

Chapter~\ref{ch:architecture} covered detection. The client-side rule, once more,
because it is the thing people get wrong:

\begin{py}
def probe_era(self) -> str:
    """Decide whether this server is `modern` or `legacy`, once."""
    try:
        self.discover()
        return "modern"
    except errors.McpError as exc:
        if exc.code in (errors.UNSUPPORTED_PROTOCOL_VERSION,
                        errors.MISSING_REQUIRED_CLIENT_CAPABILITY,
                        errors.HEADER_MISMATCH):
            return "modern"
        return "legacy"
    except Exception:
        return "legacy"
\end{py}

A recognised modern error means \emph{modern server, fix your request}. Anything
else means legacy. Cache the verdict per server process or origin, and re-probe
only if the cached assumption later fails.

\section{Health and degradation}

Servers go down. A host with four of them and one outage should be a host with
three working servers.

\begin{py}
# sketch: condensed for the page
def stats(self) -> dict:
    return {
        "servers": {
            label: {
                "healthy": b.healthy,
                "tools": len(b.tools),
                "lastError": b.last_error,
                **b.client.stats.to_json(),
            }
            for label, b in self.bindings.items()
        },
        "catalogueTools": len(self.catalogue()),
        "catalogueTokens": self.catalogue_tokens(),
    }
\end{py}

Three degradation rules:

\textbf{Drop the tools, keep the host.} An unhealthy server contributes nothing
to the catalogue. The model never sees tools it cannot call, so it never plans
around them.

\textbf{Tell the model as well as the log.} If a capability the user asked for
is
unavailable, that belongs in the context so the model can say so and stop.

\textbf{Reconnect with backoff and jitter.} A dozen hosts reconnecting to a
recovering server on a fixed one-second timer is how a recovering server stops
recovering.

That rule has two halves and people implement one of them. Backoff fixes the
\emph{rate}: doubling the delay means a client that has been failing for two
minutes is no longer asking twelve times a second. What it does not fix is the
\emph{synchronisation}, because every client that noticed the outage at the same
moment doubles in lockstep and they all arrive together anyway, just less often.

Jitter is the half that fixes the synchronisation, and full jitter is the
version worth using: sleep a uniform random amount between zero and the current ceiling. The fleet
then smears across the whole window. Adding a small wobble to the ceiling keeps
everybody clustered at its edge, which is most of the problem still.

\begin{py}
def next_delay(self, random_fraction: float) -> float:
    """The delay before attempt N. `random_fraction` is uniform in [0, 1).

    Passed in rather than drawn here so the behaviour is testable, and so a
    caller can substitute a better source.
    """
    ceiling = min(self.max_seconds, self.base_seconds * (2 ** self.attempt))
    self.attempt += 1
    return ceiling * random_fraction
\end{py}

And the part everybody forgets, which turns a blip into an outage:

\begin{py}
def reset(self) -> None:
    """Call on a successful connection, or the next outage starts at the
    ceiling and a one-second blip costs a minute of downtime."""
    self.attempt = 0
\end{py}

Without that reset, a client that had a bad afternoon three hours ago is still
carrying a 60-second ceiling, and the next one-second network hiccup costs it a
minute of unavailability for no reason at all.

\section{Building a minimal host}

Meridian's, in full, is about thirty lines:

\begin{py}
host = Host(
    capabilities=ClientCapabilities(
        elicitation={"form": {}, "url": {}},
        extensions={tasks_ext.EXTENSION_ID: {}},
    ),
    consent=ConsentPolicy(auto_allow_read_only=True),
    tool_token_budget=4000,
)

host.connect("risk", StreamableHttpClient("http://localhost:8931/mcp"),
             input_provider=FormPrompt())
host.connect("compliance", StreamableHttpClient("http://localhost:8932/mcp"),
             input_provider=FormPrompt())
host.connect("fraud", StreamableHttpClient("http://localhost:8933/mcp"))
host.connect("marketdata", StreamableHttpClient("http://localhost:8934/mcp"))

host.discover_all()        # parallel
host.refresh_catalogue()   # parallel

result = AgentLoop(host, model, max_steps=8,
                   token_budget=60_000,
                   cost_budget_usd=0.50).run("Full review of ACC-1042.")
\end{py}

Every argument there is a decision Chapter~\ref{ch:agent-loop} will justify. The
important thing is that they are all in one place, visible, and set to a number.

\section{Instrumenting from the start}

Retrofitting measurement is miserable, and Chapter~\ref{ch:measuring} depends on
these existing.

\begin{py}
@dataclass
class CallStats:
    """The counters Chapter 16 turns into a dashboard."""
    requests: int = 0
    mrtr_rounds: int = 0
    retries: int = 0
    errors: int = 0
    bytes_sent: int = 0
    bytes_received: int = 0
    latency_ms: list[float] = field(default_factory=list)
\end{py}

\texttt{mrtr\_rounds} deserves special mention. It is the count of extra round
trips forced by servers asking questions, and it is the single best early
indicator that a tool schema is missing a field. If it climbs, somebody is
asking at runtime for something they could have taken as a parameter.

Trace context propagation is one field, and it is what makes
Chapter~\ref{ch:measuring} possible at all:

\begin{py}
def _envelope(self, method, params=None, *, progress_token=None,
              traceparent=None):
    body = dict(params or {})
    body["_meta"] = build_request_meta(
        self.capabilities, self.info,
        protocol_version=self.protocol_version,
        progress_token=progress_token,
        traceparent=traceparent,
    )
    return jsonrpc.Request(self.ids.next(), method, body).to_json()
\end{py}

W3C trace context in \texttt{\_meta}, under the reserved \texttt{traceparent}
key. That is one line of code and it is the difference between four systems with
four sets of logs and one system you can ask a question of.

\subsection{What a \texttt{traceparent} contains}

Distributed tracing rests on one idea: every operation carries an identifier for
the whole request and an identifier for its own step, and every call it makes
passes those along. Reassembling the tree afterwards is then a matter of
grouping by the first and linking by the second.

W3C Trace Context is the interoperable spelling of that idea, and the whole
thing fits in a header-shaped string of four hyphen-separated fields:

\begin{plain}
00-0af7651916cd43dd8448eb211c80319c-00f067aa0ba902b7-01
|  |                                |                |
|  |                                |                +- flags: 01 = sampled
|  |                                +------------------ parent span, 8 bytes
|  +--------------------------------------------------- trace id, 16 bytes
+------------------------------------------------------ version
\end{plain}

The \textbf{trace id} is the same for every operation in one user request, which is what makes ``show me everything that happened when Ade asked for that
review'' a single query. The \textbf{parent span id}
names the specific operation that made this call, which is what turns a flat
list into a tree with causality in it. The \textbf{flags} byte carries the
sampling decision, made once at the top and honoured all the way down, so you
never get a trace that is half recorded.

Propagation is mechanical. The host has a current span, it writes
\texttt{traceparent} into \texttt{\_meta} naming that span as the parent, the
server reads it and starts its own span underneath, and if the server calls
another server it does the same thing again.

Three keys are reserved for this, and they are the only \texttt{\_meta} keys
exempt from the reverse-DNS prefix rule from Chapter~\ref{ch:prompts}:
\texttt{traceparent}, \texttt{tracestate} for vendor-specific data alongside it,
and \texttt{baggage} for arbitrary key-value pairs that ride along the whole
trace. The exemption exists because these names were already load-bearing in
every tracing library that exists, and renaming them would have bought nothing.

\begin{dangerbox}{\texttt{baggage} travels further than you think}
\texttt{baggage} propagates to every downstream service, including ones you do
not operate, and it is not encrypted.

Anything you put in it is visible to every server in the call tree. Tenant ids
and feature flags are reasonable. User email addresses are a data-sharing
decision you probably did not intend to make.
\end{dangerbox}

The payoff for one line of plumbing is that when a task takes 40 seconds, the question ``which of my four servers was slow'' is answered by opening a
trace. The alternative is correlating timestamps across four log files whose
clocks disagree.

\section{What to remember}

\begin{itemize}[leftmargin=1.6em, itemsep=3pt]
\item One client per server. The host owns everything else.
\item Discover and refresh in parallel, and let one dead server stay dead
without taking the host with it.
\item Namespace every tool, deterministically, and never change the separator.
\item The host is the context window's memory allocator. Set a budget, measure
it, put it on a dashboard.
\item A denied call is a tool execution error the model can adapt to, not an
exception.
\item Audit before executing.
\item Show arguments in consent prompts, allow scoped standing approvals, batch
related prompts.
\item The default input provider declines.
\item Probe the era once, cache the verdict, and read the error body rather than
the status code.
\item Instrument on day one, including \texttt{traceparent}. \texttt{mrtr\_rounds}
is your early warning for a missing schema field.
\end{itemize}

Next: the loop that ties it together, and the budgets that stop it running
forever.
